\SeventhArticle{8}
Now, to find the integral $\int Y\,\mathrm{d}t$, let us consider the individual parts of $Y$. Let us denote the product
\[
(t-a)(t-a')(t-a'')(t-a''')\ldotp\ldotp\ldotp\ldotp(t-a^{(n)})
\]
by $T$. Expanding this product gives
\[
T=t^{n+1}+\alpha t^n+\alpha't^{n-1}+\alpha''t^{n-2}+\text{etc.}+\alpha^{(n)}
\]
The numerator of the fraction by which, in the first part of $Y$, $A$ is multiplied is $=\frac{T}{t-a}$; the numerators in the following parts are likewise $\frac{T}{t-a'},\frac{T}{t-a''},\frac{T}{t-a'''}$, etc. The denominators, however, are simply the values of these numerators obtained by setting $t=a$, $t=a'$, $t=a''$, $t=a'''$, etc., respectively. Let us denote these denominators respectively by $M,M',M'',M'''$, etc., so that
\[
\begin{aligned}&Y=\frac{AT}{M(t-a)}+\frac{A'T}{M'(t-a')}+\frac{A''T}{M''(t-a'')}\\ &+\text{etc.}+\frac{A^{(n)}T}{M^{(n)}(t-a^{(n)})}\end{aligned}
\]
Since $T=0$ for $t=a$, we have the identity
\[
a^{n+1}+\alpha a^n+\alpha'a^{n-1}+\alpha''a^{n-2}+\text{etc.}+\alpha^{(n)}=0
\]
and therefore
\[
\begin{aligned}
T={}&t^{n+1}-a^{n+1}+\alpha(t^n-a^n)+\alpha'(t^{n-1}-a^{n-1})
\\ &+\alpha''(t^{n-2}-a^{n-2})+\text{etc.}\\
&+\alpha^{(n-1)}(t-a)
\end{aligned}
\]
Dividing by $t-a$ now gives
\[
\begin{aligned}
\frac{T}{t-a}={}&t^n+at^{n-1}+aat^{n-2}+a^3t^{n-3}+\text{etc.}+a^n\\[4pt]
&+\alpha t^{n-1}+\alpha at^{(n-2)}+\alpha aat^{n-3}+\text{etc.}+\alpha a^{n-1}\\[4pt]
&\hspace{1.01em}+\alpha't^{n-2}+\alpha'at^{n-3}+\text{etc.}+\alpha'a^{n-2}\\[4pt]
&\hspace{2.01em}+\alpha''t^{n-3}+\text{etc.}+\alpha''a^{n-3}\\[4pt]
&\hspace{3.01em}+\text{etc. etc.}\\[4pt]
&\hspace{4.01em}+\alpha^{(n-1)}
\end{aligned}
\]
\SourceLeaf{183}The value of this function for $t=a$ is found to be
\[
\begin{aligned}&=(n+1)a^n+n\alpha a^{n-1}+(n-1)\alpha'a^{n-2}\\ &+(n-2)\alpha''a^{n-3}
+\text{etc.}+\alpha^{(n-1)}\end{aligned}
\]
Hence $M$ is equal to the value of $\frac{\mathrm{d}T}{\mathrm{d}t}$ for $t=a$, as is also known from another consideration. Likewise, $M',M'',M'''$, etc., will be the values of $\frac{\mathrm{d}T}{\mathrm{d}t}$ for $t=a'$, $t=a''$, $t=a'''$, etc.

Furthermore, we find the value of the integral $\int\frac{T\,\mathrm{d}t}{t-a}$, from $t=0$ to $t=1$, to be
\[
\begin{aligned}
={}&\frac{1}{n+1}+\frac{a}{n}+\frac{aa}{n-1}+\frac{a^3}{n-2}+\text{etc.}+a^n\\[4pt]
&+\frac{\alpha}{n}+\frac{\alpha a}{n-1}+\frac{\alpha aa}{n-2}+\text{etc.}+\alpha a^{n-1}\\[4pt]
&\hspace{5.01em}+\frac{\alpha'}{n-1}+\frac{\alpha'a}{n-2}+\text{etc.}+\alpha'a^{n-2}\\[4pt]
&\hspace{6.01em}+\frac{\alpha''}{n-2}+\text{etc.}+\alpha''a^{n-3}\\[4pt]
&\hspace{7.01em}+\text{etc. etc.}\\[4pt]
&\hspace{8.01em}+\alpha^{(n-1)}
\end{aligned}
\]
We shall arrange these terms in the following order:
\[
\begin{aligned}
&a^n+\alpha a^{n-1}+\alpha'a^{n-2}+\alpha''a^{n-3}+\text{etc.}+\alpha^{(n-1)}\\[4pt]
&\quad+\tfrac12(a^{n-1}+\alpha a^{n-2}+\alpha'a^{n-3}+\text{etc.}+\alpha^{(n-2)})\\[4pt]
&\quad+\tfrac13(a^{n-2}+\alpha a^{n-3}+\alpha'a^{n-4}+\text{etc.}+\alpha^{(n-3)})\\[4pt]
&\quad+\tfrac14(a^{n-3}+\alpha a^{n-4}+\alpha'a^{n-5}+\text{etc.}+\alpha^{(n-4)})\\[4pt]
&\quad+\text{etc.}\\[4pt]
&\quad+\frac{1}{n-1}(aa+\alpha a+\alpha')\\[4pt]
&\quad+\frac{1}{n}(a+\alpha)\\[4pt]
&\quad+\frac{1}{n+1}
\end{aligned}
\]
But the same quantity evidently results if, in the product obtained by multiplying the function $T$ by the infinite series
\[
t^{-1}+\tfrac12t^{-2}+\tfrac13t^{-3}+\tfrac14t^{-4}\ \text{etc.}
\]
after all terms involving powers of $t$ with negative exponents have been discarded (or, more briefly, in the part of the product that is a polynomial in $t$), $t$ is replaced by $a$. Let us therefore set\HistoricalNote{19}
\[
T(t^{-1}+\tfrac12t^{-2}+\tfrac13t^{-3}+\tfrac14t^{-4}+\text{etc.})=T'+T''
\]
\SourceLeaf{184}where $T'$ is the polynomial in $t$ contained in this product, while $T''$ is the other part, namely a descending infinite series. This done, the value of the integral $\int\frac{T\,\mathrm{d}t}{t-a}$ from $t=0$ to $t=1$ will be equal to the value of the function $T'$ for $t=a$. Thus, if we denote the values taken by the function
\[
\frac{T'}{\left(\dfrac{\mathrm{d}T}{\mathrm{d}t}\right)}
\]
for $t=a$, $t=a'$, $t=a''$, $t=a'''$, etc., up to $t=a^{(n)}$, respectively by $R,R',R'',R'''\ldotp\ldotp\ldotp\ldotp R^{(n)}$, the integral $\int Y\,\mathrm{d}t$ from $t=0$ to $t=1$ becomes
\[
=RA+R'A'+R''A''+\text{etc.}+R^{(n)}A^{(n)}
\]
which, multiplied by $\Delta$, will give either the true or the approximate value of the integral $\int y\,\mathrm{d}x$ from $x=g$ to $x=g+\Delta$.

\SeventhArticle{9}
These operations are performed somewhat more easily if the indeterminate $t$ is replaced by another, $u=2t-1$. For brevity, we also write $b=2a-1$, $b'=2a'-1$, $b''=2a''-1$, etc. Let $T$, on replacing $t$ by $\tfrac12u+\tfrac12$, become $\frac{U}{2^{n+1}}$; that is, let
\[
U=(u-b)(u-b')(u-b'')\ldotp\ldotp\ldotp\ldotp(u-b^n)
\]
We shall therefore have $\frac{\mathrm{d}T}{\mathrm{d}t}=\frac{1}{2^n}\cdot\frac{\mathrm{d}U}{\mathrm{d}u}$, so that $M,M',M''$, etc., are the values of $\frac{1}{2^n}\cdot\frac{\mathrm{d}U}{\mathrm{d}u}$ obtained by setting successively $u=b$, $u=b'$, $u=b''$, etc.

Since the series $t^{-1}+\tfrac12t^{-2}+\tfrac13t^{-3}+\tfrac14t^{-4}+\text{etc.}$ is simply $\log.\frac{1}{1-t^{-1}}=\log\frac{1+u^{-1}}{1-u^{-1}}$, the substitution $t=\tfrac12u+\tfrac12$ necessarily transforms it into $2u^{-1}+\tfrac23u^{-3}+\tfrac25u^{-5}+\tfrac27u^{-7}+\text{etc.}$ Thus, if we set
\[
U(u^{-1}+\tfrac13u^{-3}+\tfrac15u^{-5}+\tfrac17u^{-7}+\text{etc.})=U'+U''
\]
where $U'$ is the polynomial in $u$ contained in this product, and $U''$ the other part, namely a descending infinite series, it is clear that
\[
T'+T''=\frac{1}{2^n}(U'+U'')
\]
But $T'$, being a polynomial in $t$, evidently must, under the substitution $t=\tfrac12u+\tfrac12$, yield a polynomial in $u$. By contrast, $T''$, which contains only negative powers of $t$, will yield only negative powers\SourceLeaf{185} of $u$ under the same substitution. Thus $U'$ will simply be $2^nT'$ transformed by this substitution, and likewise $U''$ will be obtained from $2^nT''$. It will therefore make no difference whether in $\dfrac{T'}{\left(\dfrac{\mathrm{d}T}{\mathrm{d}t}\right)}$ we substitute $t=a$, or in $\dfrac{U'}{\left(\dfrac{\mathrm{d}U}{\mathrm{d}u}\right)}$ we set $u=b$. We conclude that $R, R', R'', R'''$, etc., are also the values of the function $\dfrac{U'}{\left(\dfrac{\mathrm{d}U}{\mathrm{d}u}\right)}$ for $u=b$. $u=b',\ u=b'',\ u=b'''$, etc.

\SeventhArticle{10}
Before going further, we shall illustrate these rules by an example. Let $n=5$, and let us set $a=0,\ a'=\frac15,\ a''=\frac25,\ a'''=\frac35,\ a''''=\frac45,\ a'''''=1$. This gives
\[
T=t^6-3t^5+\tfrac{17}{5}t^4-\tfrac95t^3+\tfrac{274}{625}tt-\tfrac{24}{625}t
\]
Multiplying by $t^{-1}+\frac12t^{-2}+\frac13t^{-3}+\frac14t^{-4}+$ etc., we obtain
\[
T'=t^5-\tfrac52t^4+\tfrac{67}{30}t^3-\tfrac{17}{20}tt+\tfrac{913}{7500}t-\tfrac{19}{7500}
\]
The values of the coefficients $R, R', R'', R''', R'''', R'''''$ are therefore expressed by the rational function
\[
\frac{\begin{gathered}t^5-\frac25t^4+\frac{67}{30}t^3\\-\frac{17}{20}tt+\frac{913}{7500}t-\frac{19}{7500}\end{gathered}}
{\begin{gathered}6t^5-15t^4+\frac{68}{5}t^3\\-\frac{27}{5}tt+\frac{548}{625}t-\frac{24}{625}\end{gathered}}
\]
in which the successive values to be substituted for $t$ are $0,\ \frac15,\ \frac25,\ \frac35,\ \frac45,\ 1$. The other method is somewhat shorter and gives $b=-1,\ b'=-\frac35,\ b''=-\frac15,\ b'''=\frac15,\ b''''=\frac35,\ b'''''=1$
\[
\begin{aligned}
U&=u^6-\tfrac75u^4+\tfrac{259}{625}uu-\tfrac9{625}\\
U'&=u^5-\tfrac{16}{15}u^3+\tfrac{277}{1875}u
\end{aligned}
\]
\Needspace{10\baselineskip}from which $R.\ R',\ R''$, etc., will be the values of the rational function
\[
\frac{u^4-\frac{16}{15}uu+\frac{277}{1875}}
{6u^4-\frac{28}{5}uu+\frac{518}{625}}
\]
for $u=-1,\ u=-\frac35,\ u=-\frac15$, etc. Both methods give the same numbers that we presented in Article 4 from Harmonia Mensurarum. Furthermore, in a case such as this example, where $a, a', a''$, etc., are rational quantities, the values of the denominator $\frac{\mathrm{d}T}{\mathrm{d}t}$ are calculated more conveniently in its original form, namely $(a-a')(a-a'').(a-a''')\ldotp\ldotp\ldotp\ldotp(a-a^{(n)})$ for $t=a$, and likewise for the others. The same holds for the denominator $\frac{\mathrm{d}U}{\mathrm{d}u}$, which for $u=b$ becomes $=(b-b')(b-b'')(b-b''')\ldotp\ldotp\ldotp\ldotp(b-b^{(n)})$.
